\documentclass[10pt,a4paper]{extarticle}
\usepackage[a4paper,margin=10mm,top=9mm,bottom=11mm]{geometry}
% ===== Shared preamble: Class 12 Chemistry (project style, from physics-formula-sheet-source.tex) =====
\usepackage[T1]{fontenc}
\usepackage{lmodern}
\usepackage{amsmath,amssymb}
\usepackage{multicol}
\usepackage[most]{tcolorbox}
\usepackage{enumitem}
\usepackage{booktabs,tabularx,array}
\usepackage{fancyhdr}
\usepackage{xcolor}
\usepackage{microtype}
\usepackage[version=4]{mhchem}
\usepackage{tikz}
\usetikzlibrary{arrows.meta,calc}
\usepackage{pifont}
\usepackage{needspace}
\DeclareUnicodeCharacter{21CC}{\ensuremath{\rightleftharpoons}}
\DeclareUnicodeCharacter{2713}{\ding{51}}
\DeclareUnicodeCharacter{03B1}{\ensuremath{\alpha}}
\DeclareUnicodeCharacter{03B2}{\ensuremath{\beta}}
\DeclareUnicodeCharacter{2192}{\ensuremath{\to}}
\DeclareUnicodeCharacter{2610}{\ensuremath{\square}}
\DeclareUnicodeCharacter{21D2}{\ensuremath{\Rightarrow}}
\DeclareUnicodeCharacter{2194}{\ensuremath{\leftrightarrow}}
\DeclareUnicodeCharacter{2605}{\ding{72}}

% ---------- Colours: one per unit ----------
\definecolor{navy}{HTML}{1B2A4A}
\definecolor{uI}{HTML}{1F5FAD}    % 1 Solutions
\definecolor{uII}{HTML}{0E8A7E}   % 2 Electrochemistry
\definecolor{uIII}{HTML}{C2571A}  % 3 Chemical Kinetics
\definecolor{uIV}{HTML}{6A3FA0}   % 4 d- and f-Block
\definecolor{uV}{HTML}{B0204F}    % 5 Coordination Compounds
\definecolor{uVI}{HTML}{2E8B3A}   % 6 Haloalkanes & Haloarenes
\definecolor{uVII}{HTML}{8A6D00}  % 7 Alcohols, Phenols, Ethers
\definecolor{uVIII}{HTML}{2B6F8F} % 8 Aldehydes, Ketones, Acids
\definecolor{uIX}{HTML}{9A4A1E}   % 9 Amines
\definecolor{uX}{HTML}{5A5A5A}    % 10 Biomolecules
\definecolor{tipY}{HTML}{FFF4B8}
\definecolor{tipYframe}{HTML}{E0B400}
\definecolor{sqp}{HTML}{1B2A4A}
\colorlet{uc}{uI}

% ---------- Tag for questions in the official 2026-27 sample paper ----------
\newcommand{\tagbox}[2]{{\setlength{\fboxsep}{0.8pt}\colorbox{#1}{\textcolor{white}{\sffamily\bfseries\tiny\,#2\,}}}}
\newcommand{\SQP}{\tagbox{sqp}{SQP 26--27}}

% ---------- Unit banner & chapter bar ----------
\newcommand{\unitbanner}[3]{% colour, title, marks
  \colorlet{uc}{#1}%
  \par\vspace{2pt}\noindent
  \begin{tcolorbox}[enhanced,colback=#1,colframe=#1,arc=1.5pt,boxsep=1pt,left=4pt,right=4pt,top=1.5pt,bottom=1.5pt,before skip=2pt,after skip=2pt]
  \sffamily\bfseries\color{white}\footnotesize #2\hfill\scriptsize #3
  \end{tcolorbox}}
\newcommand{\chapbar}[1]{%
  \par\needspace{7\baselineskip}\noindent
  \begin{tcolorbox}[enhanced,colback=uc!18,colframe=uc,boxrule=0pt,leftrule=3pt,arc=0pt,boxsep=0.5pt,left=3pt,right=3pt,top=1pt,bottom=1pt,before skip=3pt,after skip=1.5pt]
  \sffamily\bfseries\color{uc!80!black}\footnotesize #1
  \end{tcolorbox}}
\newcommand{\sub}[1]{\par\needspace{4\baselineskip}\addvspace{2pt}\noindent{\sffamily\bfseries\scriptsize\color{uc}\strut\MakeUppercase{#1}}\par\vspace{0.5pt}}

% ---------- Boxes ----------
\newtcolorbox{kbox}{enhanced,colback=uc!5,colframe=uc,boxrule=0.6pt,arc=1.5pt,left=2pt,right=2pt,top=1pt,bottom=1pt,boxsep=1pt,before skip=2pt,after skip=2pt}
\newtcolorbox{tip}{enhanced,colback=tipY,colframe=tipYframe,boxrule=0.5pt,arc=1.5pt,left=2pt,right=2pt,top=0.5pt,bottom=0.5pt,boxsep=1pt,before skip=2pt,after skip=2pt,
  fontupper=\footnotesize}
\newtcolorbox{numbox}{enhanced,colback=white,colframe=uc!70,boxrule=0.5pt,arc=1.5pt,left=2pt,right=2pt,top=0.5pt,bottom=0.5pt,boxsep=1pt,before skip=2pt,after skip=2pt,
  borderline={0.4pt}{1.2pt}{uc!35,dashed}}
\newcommand{\Tip}{{\sffamily\bfseries\scriptsize TIP\;}}
\newcommand{\Num}{{\sffamily\bfseries\scriptsize\color{uc}STANDARD NUMERICALS\;}}

% ---------- Lists ----------
\setlist{nosep,leftmargin=8pt,labelsep=2pt}
\setlist[itemize]{label={\textcolor{uc}{\tiny$\blacktriangleright$}}}
\newenvironment{fl}{\begin{itemize}}{\end{itemize}}

% ---------- Mini graph helper (tikz) ----------
% \minigraph{width}{height}{x label}{y label}{drawing commands}
\newcommand{\minigraph}[5]{%
\begin{tikzpicture}[x=1cm,y=1cm,line cap=round,font=\tiny]
\draw[-{Stealth[length=3pt]},black!70] (0,0)--(#1,0) node[below left,inner sep=1pt]{#3};
\draw[-{Stealth[length=3pt]},black!70] (0,0)--(0,#2) node[above right,inner sep=1pt]{#4};
#5
\end{tikzpicture}}

\setlength{\columnsep}{9pt}
\setlength{\columnseprule}{0.3pt}
\def\columnseprulecolor{\color{black!25}}
\setlength{\parindent}{0pt}
\setlength{\abovedisplayskip}{1pt}\setlength{\belowdisplayskip}{1pt}
\setlength{\abovedisplayshortskip}{0pt}\setlength{\belowdisplayshortskip}{0pt}
\renewcommand{\arraystretch}{1.08}
\newcolumntype{Y}{>{\raggedright\arraybackslash}X}
\pagestyle{fancy}\fancyhf{}
\renewcommand{\headrulewidth}{0pt}

\fancyfoot[L]{\tiny\sffamily\color{black!55} Class 12 Chemistry · Must-Have Notes · Batch 1: Physical Chemistry (Ch 1--3) · CBSE Board 2026--27 (paper 27 Feb 2027)}
\fancyfoot[R]{\tiny\sffamily\color{black!55} Page \thepage}

% ---------- Notes-specific macros ----------
\newcommand{\st}[1]{\textcolor{tipYframe!90!black}{\ifcase#1\or\ding{72}\or\ding{72}\ding{72}\or\ding{72}\ding{72}\ding{72}\fi}}
\newcommand{\ty}[1]{\tagbox{uc!75!black}{#1}}
\newcommand{\mk}[1]{{\scriptsize\color{black!55}[#1]}}
\newenvironment{rqlist}{\begin{enumerate}[leftmargin=12pt,label={\sffamily\bfseries\scriptsize\color{uc}\arabic*.},itemsep=1pt]}{\end{enumerate}}
\newcommand{\ans}{\par\hspace*{2pt}{\color{uc}$\hookrightarrow$}\;}
\newtcolorbox{skel}[1]{enhanced,colback=white,colframe=uc,boxrule=0.5pt,arc=1.5pt,left=3pt,right=3pt,top=1pt,bottom=1pt,boxsep=1pt,before skip=2pt,after skip=3pt,
  title={\sffamily\bfseries\scriptsize #1},colbacktitle=uc!85,coltitle=white,toptitle=0.5pt,bottomtitle=0.5pt}
\newcommand{\chapterhead}[4]{% colour, title, marks, subtitle
  \newpage\colorlet{uc}{#1}%
  \begin{tcolorbox}[enhanced,colback=#1,colframe=#1,arc=2pt,left=6pt,right=6pt,top=3pt,bottom=3pt,after skip=3pt]
  \color{white}\sffamily{\large\bfseries #2}\hfill{\small\bfseries #3}\\[-1pt]{\scriptsize #4}
  \end{tcolorbox}}

\raggedbottom
\begin{document}
\small

% =====================================================================
% COVER
% =====================================================================
\begin{tcolorbox}[enhanced,colback=navy,colframe=navy,arc=2pt,left=8pt,right=8pt,top=6pt,bottom=6pt]
\color{white}\sffamily
{\LARGE\bfseries Class 12 Chemistry — Must-Have Notes}\\[3pt]
{\large Batch 1 of 3: Physical Chemistry} \hfill {\small Ch 1 Solutions · Ch 2 Electrochemistry · Ch 3 Chemical Kinetics}\\[2pt]
{\footnotesize CBSE Board 2026--27 only · Rationalised NCERT · Chemistry paper: \textbf{Saturday 27 February 2027} · Physical Chemistry = \textbf{23 of 70} theory marks}
\end{tcolorbox}

\colorlet{uc}{navy}
\begin{multicols}{2}
\sub{How to use these notes}
\begin{fl}
\item Each chapter is two pages: \textbf{most-repeated board questions} → \textbf{must-know NCERT lines} → \textbf{key formulas} → \textbf{answer skeletons} → \textbf{why/reason answers} → \textbf{quick 1-mark facts} → \textbf{case-based angles} → \textbf{mark-losing traps}.
\item Read with the NCERT book open. The formula sheet (Part 1) has the full formula boxes and graphs; these notes tell you \emph{what gets asked and how to write it}.
\item Try every ★★★ question on paper before reading its skeleton.
\item Tick the tracker on the last page after each read.
\end{fl}
\sub{Star legend (estimates, not exact counts)}
\begin{fl}
\item \st3\ asked in most years \quad \st2\ asked often \quad \st1\ asked sometimes.
\item Based on recurring patterns in CBSE board papers and sample papers of recent years. \SQP\ marks a question type that is in the official 2026--27 sample paper (checked).
\end{fl}
\sub{Question-type tags}
\ty{Df} define / state \;\ty{R} give reason \;\ty{N} numerical \;\ty{Dv} derivation \;\ty{G} graph / diagram \;\ty{Eq} write equations \;\ty{Cmp} difference / compare \;\ty{MCQ} 1-mark objective / A--R
\sub{Paper pattern (unchanged for 2026--27)}
\begin{fl}
\item 33 questions, 70 marks, 3 hours. A: 16 × 1 (12 MCQ + 4 A--R) · B: 5 × 2 · C: 7 × 3 · D: 2 × 4 case-based · E: 3 × 5. Internal choice in about a third of questions.
\item Physical Chemistry usually supplies: 4--6 MCQs, a 2-marker, two or three 3-markers (at least one numerical), and often one 5-marker (2026--27 SQP: Electrochemistry in Section E).
\end{fl}
\sub{Time per question (180 min)}
\begin{fl}
\item 1-mark: $\approx$1.5 min · 2-mark: 4 min · 3-mark: 6 min · case-based: 8 min · 5-mark: 10 min. That is $\approx$130 min of writing and $\approx$45 min to read the paper and check numericals.
\end{fl}

\columnbreak
\sub{Step-mark habits for Chemistry numericals}
\begin{kbox}\small
\textbf{1.} Write ``Given'' with units and convert first (g → kg for solvent, mL → L, $^\circ$C → K, min → s, S\,m$^{-1}$ → S\,cm$^{-1}$).\\
\textbf{2.} Write the formula on its own line — the formula line usually carries marks even if arithmetic slips.\\
\textbf{3.} Substitute with units in one line.\\
\textbf{4.} Box the answer \emph{with unit} and sensible significant figures. A missing or wrong unit commonly costs ½ mark.\\
\textbf{5.} For electrolytes, write ``assuming complete dissociation, $i$ = …''. For Nernst, write the balanced cell reaction and $n$.
\end{kbox}
\sub{Theory answers}
\begin{fl}
\item \textbf{Give reason}: one line stating the fact in chemistry terms + one line of cause. No stories.
\item \textbf{Equations}: balanced, with state symbols where natural; label anode/cathode for cells.
\item \textbf{Definitions/laws}: use NCERT wording — examiners mark key words (e.g.\ ``directly proportional to its mole fraction'').
\item \textbf{Graphs}: label both axes with quantity (and unit), mark intercept/slope.
\end{fl}
\sub{When to read this batch}
\begin{fl}
\item \textbf{Read 1 (now–Dec):} with NCERT text, intext questions and exercises.
\item \textbf{Read 2 (Jan, around pre-boards):} with PYQs and exemplar; mark weak spots in ``My notes''.
\item \textbf{Read 3 (17--26 Feb):} the ten days between the Physics paper (16 Feb) and Chemistry (27 Feb) are for Chemistry only — this batch on day 1--2, then timed sample papers and error review.
\end{fl}
\end{multicols}

\begin{tcolorbox}[enhanced,colback=navy!4,colframe=navy,boxrule=0.6pt,arc=2pt,left=4pt,right=4pt,top=2pt,bottom=2pt,
title={\sffamily\bfseries\footnotesize Physical Chemistry — 15 that keep coming back},colbacktitle=navy,coltitle=white]
\begin{multicols}{2}\footnotesize
\begin{enumerate}[leftmargin=14pt,itemsep=0.5pt]
\item Molar mass or new f.p./b.p. from $\Delta T_f$/$\Delta T_b$, often with $i$ (Ch 1) \st3
\item Raoult's law + positive/negative deviation + azeotrope type (Ch 1) \st3
\item van't Hoff factor, abnormal molar mass, acetic acid in benzene (Ch 1) \st3
\item Why $\pi$ for macromolecules; isotonic/hypertonic; reverse osmosis (Ch 1) \st2
\item Henry's law + scuba/altitude/soda; $K_H$ vs solubility (Ch 1) \st2
\item Nernst equation numerical, often with $\Delta G^\circ$ or $K$ (Ch 2) \st3
\item $\Lambda_m$ and $\kappa$ vs dilution; $\Lambda_m$--$\sqrt c$ graph (Ch 2) \st3
\item Kohlrausch's law + $\Lambda^\circ_m$ of weak acid, $\alpha$, $K_a$ (Ch 2) \st3
\item Lead storage battery / H$_2$--O$_2$ fuel cell / mercury cell equations (Ch 2) \st3
\item Corrosion mechanism and prevention (Ch 2) \st2
\item Faraday's law numerical; products of electrolysis (Ch 2) \st2
\item First-order numerical: $k$, $t_{1/2}$, time for $x$\% (Ch 3) \st3
\item $E_a$ from two temperatures or from $\log k$ vs $1/T$ (Ch 3) \st3
\item Order vs molecularity; pseudo first order (Ch 3) \st3
\item Derive first-order (or zero-order) integrated rate law; units of $k$ (Ch 3) \st3
\end{enumerate}
\end{multicols}
\end{tcolorbox}
\begin{tcolorbox}[enhanced,colback=white,colframe=navy,boxrule=0.8pt,arc=2pt,title={\sffamily\bfseries\footnotesize My dates \& targets (pre-board, practical exam, target marks) · My notes},colbacktitle=navy,coltitle=white,height fill,before skip=4pt]
\end{tcolorbox}

% =====================================================================
% CH 1 SOLUTIONS
% =====================================================================
\chapterhead{uI}{Ch 1 · Solutions}{7 marks}{Concentration · Henry \& Raoult · ideal/non-ideal · colligative properties · van't Hoff factor}
\begin{multicols}{2}\raggedcolumns
\sub{Most-repeated board questions}
\begin{rqlist}
\item \st3 \ty{N} Find molar mass of solute (or new f.p./b.p.) from $\Delta T_f$ or $\Delta T_b$; often an electrolyte with $i$. \mk{3}
\item \st3 \ty{Df}\ty{Cmp} State Raoult's law. What are positive and negative deviations? Give one example each and the type of azeotrope. \mk{2--3} \SQP
\item \st3 \ty{R}\ty{N} What is the van't Hoff factor? Why does ethanoic acid show half its molar mass in benzene? Find $i$ / \% dissociation from observed data. \mk{1--3}
\item \st2 \ty{R}\ty{N} Why is osmotic pressure preferred for molar mass of proteins/polymers? Calculate $\pi$ or $M_2$. \mk{2--3}
\item \st2 \ty{Df}\ty{R} State Henry's law; effect of $T$ on $K_H$ and solubility; two applications. \mk{1--2} \SQP
\item \st2 \ty{Df}\ty{R} Isotonic, hypertonic, hypotonic; what happens to blood cells in 0.5\% / 1.5\% NaCl; reverse osmosis and its use. \mk{1--2}
\item \st2 \ty{MCQ} Arrange 0.1 m solutions (glucose, NaCl, CaCl$_2$, Al$_2$(SO$_4$)$_3$) by f.p./b.p./$\pi$. \mk{1}
\item \st2 \ty{N} Molar mass from relative lowering of vapour pressure. \mk{2} \SQP
\item \st1 \ty{Df} Define molality / molarity / mole fraction; which is $T$-independent and why. \mk{1}
\item \st1 \ty{N} ppm of a gas, or mole fraction of a gas from Henry's law. \mk{1--2} \SQP
\item \st1 \ty{G} Draw $p$--$x$ diagram for positive/negative deviation, or $p$--$T$ diagram showing $\Delta T_b$. \mk{2}
\item \st1 \ty{N} Vapour pressure of a mixture of two volatile liquids, and vapour-phase mole fraction. \mk{2}
\end{rqlist}

\sub{Must-know lines (NCERT language)}
\begin{fl}
\item \textbf{Raoult's law} (volatile liquids): the partial vapour pressure of each component is directly proportional to its mole fraction in the solution.
\item For a non-volatile solute: the \textbf{relative lowering of vapour pressure} equals the mole fraction of the solute.
\item \textbf{Henry's law}: the partial pressure of a gas in the vapour phase is proportional to its mole fraction in the solution, $p=K_Hx$.
\item \textbf{Ideal solution}: obeys Raoult's law over the entire range of concentration; $\Delta_\text{mix}H=0$, $\Delta_\text{mix}V=0$.
\item \textbf{Azeotropes}: binary mixtures with the same composition in liquid and vapour phase that boil at a constant temperature.
\item \textbf{Colligative properties}: depend on the number of solute particles, not on their nature (RLVP, $\Delta T_b$, $\Delta T_f$, $\pi$).
\item \textbf{Osmotic pressure}: the extra pressure that must be applied to the solution to stop osmosis (flow of solvent into it through a semipermeable membrane).
\item \textbf{Reverse osmosis}: applying pressure larger than $\pi$ makes pure solvent flow out of the solution through the membrane.
\item \textbf{$K_b$ / $K_f$} (molal elevation / depression constant): elevation in b.p. / depression in f.p. for a 1 molal solution.
\item \textbf{Abnormal molar mass}: experimental molar mass higher or lower than the expected value, due to association or dissociation of solute.
\end{fl}

\sub{Key formulas}
\begin{kbox}
$\dfrac{p_1^\circ-p_1}{p_1^\circ}=i\,x_2$ \quad $\Delta T_b=iK_bm$ \quad $\Delta T_f=iK_fm$ \quad $\pi=iCRT$\\[2pt]
$M_2=\dfrac{1000\,K\,w_2}{\Delta T\,w_1}$ \qquad $M_2=\dfrac{w_2RT}{\pi V}$\\[2pt]$\alpha_\text{diss}=\dfrac{i-1}{n-1}$ \qquad $\alpha_\text{assoc}=\dfrac{1-i}{1-1/n}$\\[2pt]
$p_\text{total}=x_1p_1^\circ+x_2p_2^\circ$ \quad $y_1=\dfrac{x_1p_1^\circ}{p_\text{total}}$ \quad $p=K_Hx$
\end{kbox}

\sub{Answer skeletons}
\begin{skel}{N · New freezing point of an electrolyte solution (3 marks)}
\textit{e.g.\ 2 g NaCl (58.5) in 100 g water, $K_f=1.86$.}\\
Given: $w_2=2$ g, $w_1=100$ g, $M_2=58.5$, $K_f=1.86$ K\,kg\,mol$^{-1}$.\\
NaCl → Na$^+$ + Cl$^-$; assuming complete dissociation, $i=2$. \mk{½}\\
$m=\dfrac{2\times1000}{58.5\times100}=0.342$ mol kg$^{-1}$ \mk{½}\\
$\Delta T_f=iK_fm=2\times1.86\times0.342=1.27$ K \mk{1}\\
$T_f=273.15-1.27=\boxed{271.88\text{ K}}$ \;(or $-1.27\,^\circ$C) \mk{1}
\end{skel}
\begin{skel}{N · Degree of dissociation from observed data}
$i=\dfrac{\Delta T_\text{observed}}{\Delta T_\text{calculated (no dissociation)}}$ or $\dfrac{M_\text{calculated}}{M_\text{observed}}$ → $\alpha=\dfrac{i-1}{n-1}$ → state as \%.\\
Check: dissociation gives $i>1$; association gives $i<1$.
\end{skel}
\begin{skel}{N · Osmotic pressure}
Moles $n_2=w_2/M_2$; $V$ in L; $T$ in K; $R=0.083$ L\,bar\,K$^{-1}$mol$^{-1}$ (answer in bar) or $0.0821$ (atm).\\
$\pi=\dfrac{n_2RT}{V}$ → box with unit. For $M_2$: $M_2=\dfrac{w_2RT}{\pi V}$.
\end{skel}
\begin{skel}{Df+Cmp · Deviation question (3 marks)}
Definition (Raoult) → positive: A--B weaker, $p$ higher, $\Delta_\text{mix}H>0$, $\Delta_\text{mix}V>0$, e.g.\ ethanol + acetone, min-boiling azeotrope → negative: A--B stronger, $p$ lower, both $<0$, e.g.\ chloroform + acetone, max-boiling azeotrope. A small $p$--$x$ sketch earns credit if asked.
\end{skel}

\sub{Why/reason — write it like this}
\begin{rqlist}
\item \textbf{Aquatic animals are more comfortable in cold water.}\ans $K_H$ of O$_2$ increases with $T$, so solubility of O$_2$ falls with rise in $T$; cold water holds more dissolved O$_2$.
\item \textbf{Osmotic pressure is preferred for macromolecules.}\ans Measured at room $T$ (proteins denature on heating); uses molarity; $\pi$ is large enough to measure even for very dilute solutions, while $\Delta T_f$, $\Delta T_b$ would be too small.
\item \textbf{Chloroform + acetone shows negative deviation.}\ans H-bond forms between H of CHCl$_3$ and O of acetone; A--B attraction is stronger, escaping tendency and $p$ fall; $\Delta_\text{mix}H<0$.
\item \textbf{Ethanol + acetone shows positive deviation.}\ans Acetone molecules come between ethanol molecules and break some of their H-bonds; weaker A--B forces; $p$ rises.
\item \textbf{Ethanoic acid in benzene: molar mass $\approx$ 120.}\ans It dimerises through H-bonding, halving the number of particles ($i\approx0.5$); colligative property halves, so calculated molar mass doubles.
\item \textbf{Molality is preferred to molarity.}\ans Molality uses mass of solvent, which does not change with $T$; molarity uses volume, which does.
\item \textbf{0.1 m NaCl boils higher than 0.1 m glucose.}\ans NaCl gives two ions ($i=2$); glucose does not dissociate ($i=1$); $\Delta T_b\propto i$.
\item \textbf{Salt is spread on icy roads.}\ans It lowers the freezing point of water, so ice melts at lower temperature.
\item \textbf{Blood cells swell in water / shrink in 1.5\% NaCl.}\ans Water is hypotonic: water enters cells by osmosis. 1.5\% NaCl is hypertonic: water leaves cells.
\item \textbf{Azeotropes cannot be separated by fractional distillation.}\ans Liquid and vapour have the same composition at the azeotropic point, so distillation cannot enrich either component.
\item \textbf{Vapour pressure falls when a non-volatile solute is added.}\ans Solute particles occupy part of the surface, so fewer solvent molecules escape per unit time.
\end{rqlist}

\sub{Quick 1-mark facts (MCQ / A--R)}
\begin{fl}
\item Larger $K_H$ → lower solubility. Gas solubility falls as $T$ rises.
\item $T$-independent: molality, mole fraction, mass \%, ppm. $T$-dependent: molarity.
\item Positive deviation ↔ \textbf{minimum}-boiling azeotrope; negative ↔ \textbf{maximum}-boiling.
\item Ideal: $\Delta_\text{mix}H=0$, $\Delta_\text{mix}V=0$ (but $\Delta_\text{mix}S>0$, $\Delta_\text{mix}G<0$).
\item Same molality: lowest f.p./highest b.p./highest $\pi$ ↔ largest $i$. Al$_2$(SO$_4$)$_3$ (5) > CaCl$_2$ (3) > NaCl (2) > glucose (1).
\item $K_b$, $K_f$ depend only on the solvent; unit K\,kg\,mol$^{-1}$.
\item Normal saline: 0.9\% (m/V) NaCl, isotonic with blood. RO membrane: cellulose acetate.
\item Raoult's law is a special case of Henry's law ($K_H=p_1^\circ$).
\item Vapour above an ideal mixture is richer in the more volatile component.
\end{fl}

\sub{Case-based angles seen recently}
\begin{fl}
\item Scuba diving / mountaineering → Henry's law, bends, anoxia, He in tanks.
\item Desalination plant → reverse osmosis, membrane, pressure $>\pi$.
\item Hospital IV drip / saline → isotonic, what happens if hypertonic.
\item Car radiator antifreeze (ethylene glycol), salting roads → $\Delta T_f$ numerical.
\item Data table of $\Delta T_f$ for electrolytes → compute $i$, compare.
\end{fl}

\begin{tip}\Tip \textbf{Mark-losing traps.} (1) Using mass of \emph{solution} instead of solvent in molality. (2) Forgetting $i$ for NaCl, CaCl$_2$, K$_2$SO$_4$. (3) Giving $\Delta T_f$ when the question asks for the freezing point. (4) $V$ in mL or $T$ in $^\circ$C in $\pi=CRT$. (5) Swapping min-/max-boiling azeotropes. (6) Writing $p=x_2p^\circ$ for a non-volatile solute (correct: $p=x_1p_1^\circ$). (7) Saying ``higher $K_H$ means more soluble''.\end{tip}
\end{multicols}
\begin{tcolorbox}[enhanced,colback=white,colframe=uc,boxrule=0.6pt,arc=2pt,title={\sffamily\bfseries\footnotesize My notes — Ch 1 Solutions},colbacktitle=uc,coltitle=white,height fill,before skip=4pt]
\end{tcolorbox}

% =====================================================================
% CH 2 ELECTROCHEMISTRY
% =====================================================================
\chapterhead{uII}{Ch 2 · Electrochemistry}{9 marks}{Galvanic cells · Nernst · conductance · Kohlrausch · electrolysis · batteries \& fuel cells · corrosion}
\begin{multicols}{2}\raggedcolumns
\sub{Most-repeated board questions}
\begin{rqlist}
\item \st3 \ty{N} Nernst equation: $E_\text{cell}$, unknown concentration or $E^\circ$; often with $\Delta_rG^\circ$ or $K_c$. \mk{3--5} \SQP
\item \st3 \ty{G}\ty{R} How do $\kappa$ and $\Lambda_m$ vary with dilution? Plot $\Lambda_m$ vs $\sqrt c$ for strong and weak electrolytes and explain. \mk{2--3}
\item \st3 \ty{Df}\ty{N} State Kohlrausch's law; find $\Lambda^\circ_m$ of a weak acid / $\alpha$ / $K_a$ from ionic or salt data. \mk{3} \SQP
\item \st3 \ty{Eq} Reactions in the lead storage battery (discharge/charge), H$_2$--O$_2$ fuel cell, mercury cell, dry cell. \mk{2--3} \SQP
\item \st2 \ty{Eq}\ty{R} Corrosion of iron: reactions, why salt water speeds it, how galvanising / sacrificial protection works. \mk{2--3} \SQP
\item \st2 \ty{N} $\Lambda_m$ from $\kappa$, or cell constant from $R$ and $\kappa$; units of $\Lambda_m$. \mk{2--3} \SQP
\item \st2 \ty{N} Faraday's laws: mass deposited / time / charge to discharge one mole of an ion. \mk{1--2} \SQP
\item \st2 \ty{R} Products of electrolysis: aq.\ NaCl, molten NaCl/MgCl$_2$, dilute vs conc.\ H$_2$SO$_4$, CuSO$_4$ with Cu or Pt electrodes. \mk{1--2} \SQP
\item \st2 \ty{N} $\Delta_rG^\circ$ and $K_c$ from $E^\circ_\text{cell}$. \mk{2}
\item \st1 \ty{MCQ} From an $E^\circ$ table: strongest oxidising/reducing agent; can metal X displace Y; can a solution be stored in a metal vessel? \mk{1}
\item \st1 \ty{Cmp} Primary vs secondary battery; electrolytic vs galvanic cell; why mercury cell gives constant voltage. \mk{1--2}
\item \st1 \ty{Df} SHE; why absolute electrode potential cannot be measured; function of salt bridge. \mk{1--2}
\end{rqlist}

\sub{Must-know lines (NCERT language)}
\begin{fl}
\item \textbf{Galvanic cell}: converts the chemical energy of a spontaneous redox reaction into electrical energy. \textbf{Electrolytic cell}: uses electrical energy to drive a non-spontaneous reaction.
\item \textbf{Standard electrode potential}: electrode potential when all species are at unit concentration (1 M), gases at 1 bar, at 298 K.
\item \textbf{Cell potential / emf}: potential difference between the two electrodes when no current is drawn.
\item \textbf{Conductivity ($\kappa$)}: conductance of a solution of unit length and unit area of cross-section.
\item \textbf{Molar conductivity ($\Lambda_m$)}: conductivity of the solution divided by its molar concentration — the conducting power of all ions from one mole of electrolyte.
\item \textbf{Limiting molar conductivity ($\Lambda^\circ_m$)}: molar conductivity when concentration approaches zero (infinite dilution).
\item \textbf{Kohlrausch's law}: limiting molar conductivity of an electrolyte is the sum of the individual contributions of its anion and cation.
\item \textbf{Faraday's first law}: amount of chemical reaction at an electrode is proportional to the quantity of electricity passed. \textbf{Second law}: for the same charge, masses liberated are proportional to their chemical equivalent weights (atomic mass ÷ electrons needed).
\item \textbf{Primary battery}: reaction occurs once, cannot be recharged. \textbf{Secondary}: recharged by passing current in the opposite direction.
\item \textbf{Fuel cell}: galvanic cell that converts the energy of combustion of fuels (H$_2$, CH$_4$, CH$_3$OH) directly into electrical energy.
\item \textbf{Corrosion}: slow coating of a metal surface with oxides or other salts; an electrochemical process.
\end{fl}

\sub{Key formulas}
\begin{kbox}
$E^\circ_\text{cell}=E^\circ_\text{c}-E^\circ_\text{a}$ \qquad $E=E^\circ-\dfrac{0.059}{n}\log Q$\\[2pt]$\Delta_rG^\circ=-nFE^\circ$ \qquad $\log K_c=\dfrac{nE^\circ}{0.059}$\\[2pt]
$\kappa=\dfrac{G^*}{R}$ \qquad $\Lambda_m=\dfrac{1000\,\kappa}{M}$ \qquad $\Lambda_m=\Lambda^\circ_m-A\sqrt c$\\[2pt]$\Lambda^\circ_m=\nu_+\lambda^\circ_++\nu_-\lambda^\circ_-$ \qquad 
$\alpha=\dfrac{\Lambda_m}{\Lambda^\circ_m}$\\[2pt]$K_a=\dfrac{c\alpha^2}{1-\alpha}$ \quad $m=\dfrac{M\,I\,t}{nF}$
\end{kbox}

\sub{Answer skeletons}
\begin{skel}{N · Nernst for a cell (3 marks)}
\textit{e.g.\ \ce{Zn|Zn^2+(0.1 M)||Cu^2+(0.01 M)|Cu}, $E^\circ=1.10$ V.}\\
Cell reaction: \ce{Zn + Cu^2+ -> Zn^2+ + Cu}, $n=2$ \mk{½}\\
$E=E^\circ-\dfrac{0.059}{2}\log\dfrac{[\ce{Zn^2+}]}{[\ce{Cu^2+}]}$ \mk{1}\\
$=1.10-0.0295\log\dfrac{0.1}{0.01}=1.10-0.0295$ \mk{1}\\
$=\boxed{1.07\text{ V}}$ \mk{½}
\end{skel}
\begin{skel}{N · $\Lambda_m$ from resistance (3 marks)}
Cell constant $G^*=R\kappa$ (from standard KCl) → $\kappa_\text{new}=G^*/R_\text{new}$ → convert to S\,cm$^{-1}$ → $\Lambda_m=1000\kappa/M$ → unit S\,cm$^2$mol$^{-1}$.
\end{skel}
\begin{skel}{N · Kohlrausch for a weak acid (3 marks)}
Write the relation with ions: $\Lambda^\circ_{\ce{HA}}=\lambda^\circ_{\ce{H+}}+\lambda^\circ_{\ce{A-}}$, or with salts: $\Lambda^\circ_{\ce{HA}}=\Lambda^\circ_{\ce{NaA}}+\Lambda^\circ_{\ce{HCl}}-\Lambda^\circ_{\ce{NaCl}}$ \mk{1} → substitute \mk{½} → $\alpha=\Lambda_m/\Lambda^\circ_m$ \mk{½} → $K_a=c\alpha^2/(1-\alpha)$ \mk{1}.
\end{skel}
\begin{skel}{N · Faraday (2 marks)}
$Q=It$ ($t$ in s) → moles e$^-$ $=Q/F$ → moles metal $=$ moles e$^-$ ÷ $n$ → mass $=$ moles × $M$.
\end{skel}
\begin{skel}{Eq · Lead storage battery (3 marks)}
Anode: \ce{Pb + SO4^2- -> PbSO4 + 2e-}\\
Cathode: \ce{PbO2 + SO4^2- + 4H+ + 2e- -> PbSO4 + 2H2O}\\
Overall: \ce{Pb + PbO2 + 2H2SO4 -> 2PbSO4 + 2H2O}; on charging the reaction is reversed.
\end{skel}

\sub{Why/reason — write it like this}
\begin{rqlist}
\item \textbf{$\kappa$ decreases but $\Lambda_m$ increases on dilution.}\ans Fewer ions per unit volume, so $\kappa$ falls; $\Lambda_m$ counts ions from one mole in the whole volume, and these ions move more freely (strong) or are more numerous because $\alpha$ rises (weak).
\item \textbf{$\Lambda^\circ_m$ of a weak electrolyte cannot be found by extrapolation.}\ans $\Lambda_m$ rises very steeply at low $c$, so the curve cannot be extended to $c=0$; Kohlrausch's law is used instead.
\item \textbf{Mercury cell gives a constant voltage.}\ans The overall reaction \ce{Zn(Hg) + HgO -> ZnO + Hg} involves no ions whose concentration changes during use.
\item \textbf{Aqueous NaCl gives H$_2$, not Na, at the cathode.}\ans $E^\circ$(\ce{H2O/H2}) $=-0.83$ V is higher than $E^\circ$(\ce{Na+/Na}) $=-2.71$ V, so water is reduced. Cl$_2$ forms at the anode instead of O$_2$ because of the overpotential of oxygen.
\item \textbf{Iron rusts faster in sea water.}\ans Dissolved salts increase the conductivity of the water film, speeding up the electrochemical reactions.
\item \textbf{Galvanised iron is protected even when scratched.}\ans Zn has a lower reduction potential than Fe, so Zn is oxidised in preference (sacrificial anode).
\item \textbf{CuSO$_4$ solution cannot be stored in a zinc pot.}\ans $E^\circ$(\ce{Zn^2+/Zn}) $<$ $E^\circ$(\ce{Cu^2+/Cu}); Zn displaces Cu: \ce{Zn + Cu^2+ -> Zn^2+ + Cu}.
\item \textbf{AC is used to measure conductance.}\ans DC would cause electrolysis and change the concentration at the electrodes.
\item \textbf{Absolute electrode potential cannot be measured.}\ans A half-cell cannot work alone; potentials are measured against SHE (taken as 0 V).
\item \textbf{Electrolytic conductance rises with $T$.}\ans Ions move faster (lower viscosity, less solvation); for weak electrolytes dissociation also increases.
\item \textbf{Fuel cells are preferred.}\ans About 70\% efficient, pollution-free (product water), supply continuous as long as reactants are fed.
\item \textbf{Acid density in a lead battery falls on discharge.}\ans H$_2$SO$_4$ is consumed and water is formed.
\end{rqlist}

\sub{Quick 1-mark facts (MCQ / A--R)}
\begin{fl}
\item Units: $\kappa$ S\,cm$^{-1}$ (S\,m$^{-1}$), $\Lambda_m$ S\,cm$^2$mol$^{-1}$ (S\,m$^2$mol$^{-1}$), $G^*$ cm$^{-1}$, $F$ C\,mol$^{-1}$. \SQP
\item $E^\circ$ is intensive; $\Delta G$ is extensive. At equilibrium $E_\text{cell}=0$, $Q=K$.
\item Galvanic: anode $-$, cathode $+$. Electrolytic: anode $+$, cathode $-$. Oxidation is always at the anode.
\item Strongest oxidising agent: F$_2$; strongest reducing agent (aqueous): Li.
\item 1 F deposits 1 mol of a univalent ion: 108 g Ag, 31.75 g Cu, 9 g Al. Charge to get 1 mol Al from Al$^{3+}$ (or from M$_2$O$_3$): $3F$. \SQP
\item At the same concentration, HCl conducts best (very high $\lambda^\circ$ of H$^+$); CH$_3$COOH conducts poorly (weak). \SQP
\item Mercury cell 1.35 V; dry cell $\approx$1.5 V; Daniell 1.10 V; corrosion cell 1.67 V.
\item Lead battery electrolyte: 38\% H$_2$SO$_4$. Fuel cell electrolyte (NCERT): conc.\ NaOH. \SQP
\item Rust: \ce{Fe2O3.xH2O}. Sacrificial metals: Zn, Mg.
\end{fl}

\sub{Case-based angles seen recently}
\begin{fl}
\item Table of $\Lambda_m$ at several $c$ → identify strong/weak, find $\Lambda^\circ_m$, comment on the graph.
\item Hydrogen buses / space missions → fuel cell reactions, advantages.
\item Car battery / inverter → lead storage battery, charging, acid density.
\item Ships, pipelines, bridges → corrosion, sacrificial anodes, galvanising.
\item Electroplating jewellery → Faraday numerical.
\end{fl}

\begin{tip}\Tip \textbf{Mark-losing traps.} (1) Multiplying $E^\circ$ by a coefficient. (2) Writing $Q$ upside down or including solids. (3) Wrong $n$ (take it from the \emph{balanced overall} equation). (4) $\kappa$ in S\,m$^{-1}$ used directly in $1000\kappa/M$ (1 S\,m$^{-1}$ = 0.01 S\,cm$^{-1}$). (5) Time in minutes in $Q=It$. (6) Forgetting coefficients in Kohlrausch, e.g.\ $\Lambda^\circ_{\ce{Al2(SO4)3}}=2\lambda^\circ_{\ce{Al^3+}}+3\lambda^\circ_{\ce{SO4^2-}}$. (7) Writing fuel-cell reactions in acid when NCERT's alkaline version is expected.\end{tip}
\end{multicols}
\begin{tcolorbox}[enhanced,colback=white,colframe=uc,boxrule=0.6pt,arc=2pt,title={\sffamily\bfseries\footnotesize My notes — Ch 2 Electrochemistry},colbacktitle=uc,coltitle=white,height fill,before skip=4pt]
\end{tcolorbox}

% =====================================================================
% CH 3 CHEMICAL KINETICS
% =====================================================================
\chapterhead{uIII}{Ch 3 · Chemical Kinetics}{7 marks}{Rate · rate law · order \& molecularity · integrated equations · half-life · Arrhenius · collision theory}
\begin{multicols}{2}\raggedcolumns
\sub{Most-repeated board questions}
\begin{rqlist}
\item \st3 \ty{N} First-order numerical: $k$ from data, $t_{1/2}$, time for 75/90/99\% completion; ``show that $t_{99.9\%}=10\,t_{1/2}$''. \mk{2--3} \SQP
\item \st3 \ty{N}\ty{G} $E_a$ from $k$ at two temperatures or from a $\log k$ vs $1/T$ plot; effect of catalyst on $E_a$. \mk{3} \SQP
\item \st3 \ty{Cmp}\ty{Df} Order vs molecularity (two differences); pseudo first-order reaction with example. \mk{2} \SQP
\item \st3 \ty{Dv} Derive the integrated rate equation for a first-order (or zero-order) reaction and its half-life. \mk{3}
\item \st2 \ty{MCQ} Units of $k$ for a given order, or order from the units. \mk{1}
\item \st2 \ty{N} Order from initial-rate data; write the rate law and find $k$. \mk{3}
\item \st2 \ty{G} Identify order from a graph ($[R]$--$t$, $\ln[R]$--$t$, rate--$[R]$, $t_{1/2}$--$[R]_0$); meaning of slope/intercept. \mk{1--2}
\item \st2 \ty{MCQ} Rate in terms of each species for a given equation, e.g.\ \ce{2N2O5 -> 4NO2 + O2}. \mk{1--2}
\item \st2 \ty{G}\ty{R} Effect of temperature (Maxwell--Boltzmann curves) and catalyst (energy profile). \mk{2}
\item \st1 \ty{Df} Collision theory: effective collisions, orientation, steric factor $P$. \mk{1--2}
\item \st1 \ty{N} Zero-order: $k$ and $t_{1/2}$; rate stays constant. \mk{1--2} \SQP
\item \st1 \ty{N} First-order gas-phase reaction from total pressure. \mk{3}
\end{rqlist}

\sub{Must-know lines (NCERT language)}
\begin{fl}
\item \textbf{Rate law}: rate expressed in terms of molar concentrations of reactants, each raised to a power that may or may not equal its stoichiometric coefficient.
\item \textbf{Order}: sum of the powers of the concentration terms in the rate law; found by experiment; can be 0, fractional or integer.
\item \textbf{Molecularity}: number of reacting species (atoms, ions, molecules) that must collide simultaneously in an elementary reaction; 1, 2 or 3.
\item \textbf{Rate constant}: rate of reaction when the concentration of each reactant is unity.
\item \textbf{Half-life}: time in which the concentration of a reactant falls to half its initial value.
\item \textbf{Pseudo first order}: a reaction of higher order that behaves as first order because one reactant is in large excess.
\item \textbf{Activation energy}: extra energy that reactant molecules must absorb to form the activated complex.
\item \textbf{Catalyst}: increases rate without being consumed, by providing an alternative path of lower $E_a$; does not alter $\Delta G$ or the equilibrium constant; catalyses only spontaneous reactions.
\item \textbf{Effective collisions}: collisions with energy $\ge$ threshold energy \emph{and} proper orientation.
\end{fl}

\sub{Key formulas}
\begin{kbox}
Zero: $k=\dfrac{[R]_0-[R]}{t}$, \;$t_{1/2}=\dfrac{[R]_0}{2k}$\\[2pt]First: $k=\dfrac{2.303}{t}\log\dfrac{[R]_0}{[R]}$, \;$t_{1/2}=\dfrac{0.693}{k}$\\[2pt]
Gas: $k=\dfrac{2.303}{t}\log\dfrac{p_i}{2p_i-p_t}$ \qquad $k=Ae^{-E_a/RT}$\\[2pt]$\log\dfrac{k_2}{k_1}=\dfrac{E_a}{2.303R}\cdot\dfrac{T_2-T_1}{T_1T_2}$\\[2pt]
Units of $k$: $(\text{mol L}^{-1})^{1-n}\text{s}^{-1}$ \quad Rate $=PZ_{AB}e^{-E_a/RT}$
\end{kbox}

\sub{Answer skeletons}
\begin{skel}{Dv · First-order integrated rate law (3 marks)}
\ce{R -> P}; \;rate $=-\dfrac{d[R]}{dt}=k[R]$ \mk{½}\\
$\dfrac{d[R]}{[R]}=-k\,dt$ → integrate: $\ln[R]=-kt+I$ \mk{½}\\
At $t=0$, $[R]=[R]_0$ ⇒ $I=\ln[R]_0$ \mk{½}\\
$\ln[R]=-kt+\ln[R]_0$ ⇒ $k=\dfrac1t\ln\dfrac{[R]_0}{[R]}=\boxed{\dfrac{2.303}{t}\log\dfrac{[R]_0}{[R]}}$ \mk{½}\\
At $t_{1/2}$, $[R]=[R]_0/2$: $t_{1/2}=\dfrac{2.303\log2}{k}=\boxed{\dfrac{0.693}{k}}$ \mk{1}
\end{skel}
\begin{skel}{Dv · Zero order (2--3 marks)}
$-\dfrac{d[R]}{dt}=k$ → $d[R]=-k\,dt$ → $[R]=-kt+I$; $I=[R]_0$ → $[R]=[R]_0-kt$, $k=\dfrac{[R]_0-[R]}{t}$; $t_{1/2}=\dfrac{[R]_0}{2k}$.
\end{skel}
\begin{skel}{N · Time for $x$\% completion (2 marks)}
Take $[R]_0=100$, $[R]=100-x$. \;$t=\dfrac{2.303}{k}\log\dfrac{100}{100-x}$. Find $k$ from $t_{1/2}$ if needed: $k=0.693/t_{1/2}$.\\
Shortcuts to verify: 75\% → $2t_{1/2}$; 87.5\% → $3t_{1/2}$; 99.9\% → $\approx10t_{1/2}$.
\end{skel}
\begin{skel}{N · $E_a$ from two temperatures (3 marks)}
Write the two-point Arrhenius form \mk{1} → substitute with $R=8.314$, $T$ in K \mk{1} → $E_a$ in J mol$^{-1}$, convert to kJ mol$^{-1}$ \mk{1}.
\end{skel}
\begin{skel}{N · Order from initial rates (3 marks)}
Compare two runs where only [A] changes: $\dfrac{r_2}{r_1}=\left(\dfrac{[A]_2}{[A]_1}\right)^x$ → $x$. Same for $y$. Write rate $=k[A]^x[B]^y$, overall order, then $k$ with units from one run.
\end{skel}

\sub{Why/reason — write it like this}
\begin{rqlist}
\item \textbf{Rate increases with temperature.}\ans More molecules have energy $\ge E_a$; for a 10 K rise the fraction (area beyond $E_a$) roughly doubles, so $k$ nearly doubles.
\item \textbf{Order can be fractional; molecularity cannot.}\ans Order is an experimental quantity from the rate law; molecularity counts molecules colliding in one step, which must be a whole number.
\item \textbf{Molecularity above 3 is not observed.}\ans The chance of more than three molecules colliding at the same instant with proper energy and orientation is very small.
\item \textbf{Hydrolysis of ethyl acetate in water is pseudo first order.}\ans Water is in large excess, so [H$_2$O] stays practically constant; rate $=k'[\text{ester}]$.
\item \textbf{A catalyst does not change the equilibrium constant.}\ans It lowers $E_a$ of the forward and backward reactions by the same amount, so both speed up equally; $\Delta G$ is unchanged.
\item \textbf{Not every collision gives products.}\ans Only collisions with sufficient energy and proper orientation are effective.
\item \textbf{$t_{1/2}$ of a first-order reaction is constant.}\ans $t_{1/2}=0.693/k$ contains no concentration term.
\item \textbf{Rate decreases as the reaction proceeds.}\ans Concentration of reactants decreases.
\item \textbf{Molecularity has no meaning for a complex reaction.}\ans A complex reaction has several steps; molecularity is defined for each elementary step; the slowest step fixes the rate.
\end{rqlist}

\sub{Quick 1-mark facts (MCQ / A--R)}
\begin{fl}
\item Rate always has unit mol L$^{-1}$s$^{-1}$; $k$'s unit depends on order: s$^{-1}$ → first; L mol$^{-1}$s$^{-1}$ → second; mol L$^{-1}$s$^{-1}$ → zero.
\item \ce{2N2O5 -> 4NO2 + O2}: $-\frac12\frac{d[\ce{N2O5}]}{dt}=\frac14\frac{d[\ce{NO2}]}{dt}=\frac{d[\ce{O2}]}{dt}$.
\item Rate $=k[A]^{1/2}[B]^{3/2}$ → overall order 2.
\item $k$ depends on $T$ and catalyst, not on concentration.
\item Zero order: rate constant; $t_{1/2}\propto[R]_0$. First order: $t_{1/2}$ independent of $[R]_0$.
\item $\ln k$ vs $1/T$: slope $-E_a/R$; $\log k$ vs $1/T$: slope $-E_a/2.303R$; intercept $\ln A$ / $\log A$. As $T\to\infty$, $k\to A$.
\item A catalyst lowers $E_a$, does not change $\Delta H$, $\Delta G$ or $K$.
\item Radioactive decay is first order (half-life questions on $^{14}$C use the first-order formula).
\item Molecularity is meaningful only for elementary reactions; it is never zero.
\end{fl}

\sub{Case-based angles seen recently}
\begin{fl}
\item Concentration--time table → test constant $t_{1/2}$ or constant $k$ to decide order.
\item Food spoilage / refrigeration, cooking at altitude → temperature effect, Arrhenius.
\item Medicine shelf-life / drug in blood → first-order $t_{1/2}$, time to fall to $x$\%.
\item Enzymes, catalytic converters → catalyst and energy profile.
\item Hydrolysis of acetyl chloride / ester in water → pseudo first order. \SQP
\end{fl}

\begin{tip}\Tip \textbf{Mark-losing traps.} (1) Taking order from the balanced equation instead of data. (2) Mixing $\ln$ and $\log$ (forgetting 2.303). (3) Time units: $k$ in min$^{-1}$ gives $t$ in min. (4) Arrhenius: $E_a$ in kJ with $R=8.314$ J, or $T$ in $^\circ$C. (5) Rate of reaction vs rate of disappearance of one species (stoichiometric factors). (6) ``Catalyst increases $E_a$'' or ``shifts equilibrium''. (7) Assigning molecularity to a complex reaction.\end{tip}
\end{multicols}
\begin{tcolorbox}[enhanced,colback=white,colframe=uc,boxrule=0.6pt,arc=2pt,title={\sffamily\bfseries\footnotesize My notes — Ch 3 Chemical Kinetics},colbacktitle=uc,coltitle=white,height fill,before skip=4pt]
\end{tcolorbox}

% =====================================================================
% PROGRESS TRACKER
% =====================================================================
\newpage
\colorlet{uc}{navy}
\begin{tcolorbox}[enhanced,colback=navy,colframe=navy,arc=2pt,left=6pt,right=6pt,top=3pt,bottom=3pt]
\color{white}\sffamily{\large\bfseries Progress tracker — Class 12 Chemistry}\hfill{\small tick ☐ as you go}\\[-1pt]{\scriptsize Notes for Ch 4--10 come in Batches 2 (Inorganic) and 3 (Organic); the rows are here so one page tracks the whole subject.}
\end{tcolorbox}
{\small
\renewcommand{\arraystretch}{1.55}
\begin{tabularx}{\textwidth}{@{}l@{\ }>{\raggedright\arraybackslash}p{4.7cm}*{8}{>{\centering\arraybackslash}X}@{}}
\toprule
 & \textbf{Chapter (marks)} & \textbf{Read 1} & \textbf{Read 2} & \textbf{Read 3} & \textbf{NCERT Qs} & \textbf{Exemplar} & \textbf{PYQs} & \textbf{Sheet drilled} & \textbf{Conf.\ (1--5)}\\
\midrule
\textcolor{uI}{\textbf{1}} & Solutions (7) & ☐ & ☐ & ☐ & ☐ & ☐ & ☐ & ☐ & \underline{\hspace{8mm}}\\
\textcolor{uII}{\textbf{2}} & Electrochemistry (9) & ☐ & ☐ & ☐ & ☐ & ☐ & ☐ & ☐ & \underline{\hspace{8mm}}\\
\textcolor{uIII}{\textbf{3}} & Chemical Kinetics (7) & ☐ & ☐ & ☐ & ☐ & ☐ & ☐ & ☐ & \underline{\hspace{8mm}}\\
\textcolor{uIV}{\textbf{4}} & d- and f-Block Elements (7) & ☐ & ☐ & ☐ & ☐ & ☐ & ☐ & ☐ & \underline{\hspace{8mm}}\\
\textcolor{uV}{\textbf{5}} & Coordination Compounds (7) & ☐ & ☐ & ☐ & ☐ & ☐ & ☐ & ☐ & \underline{\hspace{8mm}}\\
\textcolor{uVI}{\textbf{6}} & Haloalkanes \& Haloarenes (6) & ☐ & ☐ & ☐ & ☐ & ☐ & ☐ & ☐ & \underline{\hspace{8mm}}\\
\textcolor{uVII}{\textbf{7}} & Alcohols, Phenols \& Ethers (6) & ☐ & ☐ & ☐ & ☐ & ☐ & ☐ & ☐ & \underline{\hspace{8mm}}\\
\textcolor{uVIII}{\textbf{8}} & Aldehydes, Ketones \& Acids (8) & ☐ & ☐ & ☐ & ☐ & ☐ & ☐ & ☐ & \underline{\hspace{8mm}}\\
\textcolor{uIX}{\textbf{9}} & Amines (6) & ☐ & ☐ & ☐ & ☐ & ☐ & ☐ & ☐ & \underline{\hspace{8mm}}\\
\textcolor{uX}{\textbf{10}} & Biomolecules (7) & ☐ & ☐ & ☐ & ☐ & ☐ & ☐ & ☐ & \underline{\hspace{8mm}}\\
\bottomrule
\end{tabularx}}

\vspace{4pt}
{\small\textbf{Practicals (30):} Volumetric~☐ \quad Salt analysis~☐ \quad Content-based experiments~☐ \quad Project~☐ \quad Record complete~☐ \quad Viva prepared~☐\\[3pt]
\textbf{Timed sample papers (3 h, error review after each):} SQP 2026--27~☐ \quad Paper 2~☐ \quad Paper 3~☐ \quad Paper 4~☐ \quad Paper 5~☐ \quad Paper 6~☐}
\vspace{2pt}
\begin{tcolorbox}[enhanced,colback=tipY,colframe=tipYframe,boxrule=0.5pt,arc=1.5pt,left=4pt,right=4pt,top=2pt,bottom=2pt]
\small\textbf{Error log habit.} After every test, write each lost mark in ``My notes'' under one of: \emph{concept gap} · \emph{formula/unit slip} · \emph{NCERT line not known} · \emph{time}. Re-read only the matching section of these notes.
\end{tcolorbox}

\begin{tcolorbox}[enhanced,colback=white,colframe=navy,boxrule=0.8pt,arc=2pt,title={\sffamily\bfseries\footnotesize My notes},colbacktitle=navy,coltitle=white,height fill,before skip=4pt]
\end{tcolorbox}
\end{document}
